
\slajdid{05-s049}
\begin{frame}[label=05-s049]
\gusttekst
\begin{center}ODSECANJE REZULTATA MNOŽENJA\end{center}
\[a_{n-1}a_{n-2}\ldots a_1a_0\times b_{n-1}b_{n-2}\ldots b_1b_0=c_{2n-1}c_{2n-2}\ldots c_n\textcolor{DEcrvena}{c_{n-1}c_{n-2}\ldots c_1c_0}\]
Ako rezultat množenja smeštamo takođe u $n$ bitni registar kod celobrojnog množenja mora da se vodi računa da ne dođe do prekoračenja opsega, odnosno da rezultat može da stane u $n$ bitni registar
\begin{align*}
c_{2n-1}c_{2n-2}\ldots c_n c_{n-1}&=00\ldots00&&\text{Za pozitivan rezultat, označeno množenje}\\[3mm]
c_{2n-1}c_{2n-2}\ldots c_n c_{n-1}&=11\ldots11&&\text{Za negativan rezultat, označeno množenje}\\[3mm]
c_{2n-1}c_{2n-2}\ldots c_n&=00\ldots00&&\text{Neoznačeno množenje}
\end{align*}
Greška koju bi napravili eventualnim prekoračenjem opsega bila bi „ogromna“.\\Videli smo to kod sabiranja.
\end{frame}
